\block{\textsc{Most used method}}{
    A \textbf{neural network} is a computational model inspired by the human brain. It consists of layers of interconnected nodes called neurons. At a minimum, a network has an input layer and an output layer. In addition to that, we can add one or more hidden layers between them to complexify the function that the network can approximate. Information from the input layer propagates through the hidden layers to the output layer. Each neuron computes a weighted sum of its inputs and applies an activation function to this sum. Neural networks can learn complex patterns from data, making them very powerful tools.\\
    \begin{center}
    \includegraphics[width=20cm]{images/OIP.jpg}    
    \end{center}
    
    \vspace{0.5cm}
    \textbf{Recurrent neural network} (RNNs) are a class of neural networks designed to handle sequential data, where the output at each step depends on previous computations. They maintain a hidden state that captures information from prior inputs, making them effective for time series forecasting, natural language processing, and sequence prediction. Unlike feedforward networks, RNNs can model temporal dependencies. However, standard RNNs may struggle with long-term dependencies due to vanishing or exploding gradients. Advanced variants like Long Short-Term Memory (LSTM) address these limitations by improving the network’s ability to learn from distant past events.
    \begin{center}
        \includegraphics[width= 15 cm]{images/enterpriseai-recurrent_neural_network-f_mobile.png}
    \end{center}
    \vspace{0.5cm}
    \textbf{Support Vector Machines} (SVM) are supervised learning models widely used for both regression and classification tasks in machine learning. In the context of regression, known as Support Vector Regression (SVR), the algorithm seeks to fit a function within a margin of tolerance … margin of tolerance $\epsilon$ around … around the predicted values. Data points that fall inside this margin do not contribute to the error, whereas the optimization focuses on minimizing the error of points lying outside the margin. The estimated function can model nonlinear relationships through the use of nonlinear kernel functions. Common examples of nonlinear kernels include the polynomial kernel and the Gaussian (RBF) kernel. With an appropriate choice of kernel, SVR can act as a universal approximator, capable of modeling a wide range of complex functions.
    \\

    \begin{center}
        \includegraphics[width = 15 cm]{images/support-vector-regression-svr1-l.jpg} 
    \end{center}
}      